\documentclass[11pt,a4paper]{article}
\usepackage[margin=18mm,headheight=15pt]{geometry}
\usepackage[T1]{fontenc}\usepackage{lmodern,amsmath,amssymb,booktabs,array,enumitem,fancyhdr,xcolor}
\usepackage[american]{circuitikz}\usepackage[hidelinks]{hyperref}
\definecolor{ink}{HTML}{17324D}\definecolor{pale}{HTML}{EDF3F8}
\pagestyle{fancy}\fancyhf{}\lhead{\small BJT practice workbook}\rhead{\small $R_{\rm in}$, $R_{\rm out}$ and $A_v$}\cfoot{\small\thepage}
\setlength{\emergencystretch}{3em}\setlength{\parindent}{0pt}\setlength{\parskip}{6pt}
\setlist[enumerate]{leftmargin=*,itemsep=4pt,topsep=3pt}
\newcommand{\K}{\,\mathrm{k}\Omega}\newcommand{\Oh}{\,\Omega}\newcommand{\mS}{\,\mathrm{mS}}
\newcommand{\head}[1]{\par{\Large\bfseries\color{ink}#1}\par\vspace{3pt}\hrule\vspace{7pt}}
\newcommand{\ex}[2]{\par\vspace{5pt}{\large\bfseries\color{ink}Exercise #1 --- #2}\label{ex:#1}\par}
\newcommand{\sol}[2]{\par{\large\bfseries\color{ink}Solution #1 --- #2}\label{sol:#1}\par}
\newcommand{\work}{\vfill\textcolor{gray}{\rule{\linewidth}{.2pt}}\par\vspace{10pt}\textcolor{gray}{\rule{\linewidth}{.2pt}}\par\vspace{10pt}}
\begin{document}

{\Huge\bfseries\color{ink}BJT amplifiers}\par
{\LARGE Input resistance, output resistance, and voltage gain}\par
\textbf{24 progressive exercises with diagrams and worked solutions}

Built around the BJT amplifier topics in your Lectures 5--8 and Sheets 6--7. All exercises are new. Start with the first three pages, then solve the questions before opening the solutions.

\head{The five-step routine}
\begin{enumerate}
\item \textbf{Find the DC bias first, if it is not given.} Capacitors are open for DC. Compute $I_C$ and verify the region of operation. Then use $g_m=I_C/V_T$ and $r_\pi=\beta/g_m$.
\item \textbf{Redraw for midband AC.} Ground ideal DC voltage supplies; open ideal DC current sources. Short sufficiently large coupling/bypass capacitors. Keep every resistor that still connects two different AC nodes.
\item \textbf{Name the ports.} Here $v_i$ means the actual voltage at the amplifier input terminal. The source voltage $v_s$ is usually different because $R_s$ and $R_{\rm in}$ form a divider.
\item \textbf{Find loaded voltage gain.} Include $R_L$, or the next stage's input resistance, in the output load. Find $A_v=v_o/v_i$ first, then $A_s=v_o/v_s=A_vR_{\rm in}/(R_s+R_{\rm in})$ when the input divider applies.
\item \textbf{Find output resistance independently.} Remove the external load, set independent signal sources to zero, retain their internal resistances, and apply a test voltage $v_x$. Keep the controlled source: $R_{\rm out}=v_x/i_x$.
\end{enumerate}
\textbf{What setting a source to zero means:} an ideal voltage source becomes a short; an ideal current source becomes an open. A voltage source with series resistance leaves that resistance in the circuit.

\head{A port checklist}
\begin{center}\begin{tabular}{p{3.2cm}p{11.6cm}}\toprule
Quantity & What you must include\\\midrule
$R_{\rm in}$ & Everything connected inside the input port, including bias resistors. Exclude the source's own series resistance.\\
$A_v=v_o/v_i$ & The actual output load. Specify whether $v_i$ is at base or emitter.\\
$A_s=v_o/v_s$ & Source attenuation as well as loaded stage gain.\\
$R_{\rm out}$ & Source resistance after zeroing its voltage. Remove only the external load, not an internal emitter or collector resistor.\\\bottomrule
\end{tabular}\end{center}
\textbf{Conventions throughout.} $V_T=25$ mV; transistors are forward-active; $r_o=\infty$ unless stated. $R_B$ means the equivalent AC bias resistance, usually $R_1\parallel R_2$. Figures are midband AC circuits with bias sources omitted, except the explicitly labeled DC exercise; bypass capacitors are shown where their connections are being tested. All numerical gains are V/V. Use finite $\beta$ unless an approximation is requested.

\textbf{Suggested route.} First solve 1--6, then 9--14, then 7--8 and 15--24. Allow about 10--15 minutes per exercise, longer for cascades. Solutions begin at page \pageref{solutions}.
\newpage
\head{A compact map of the formulas}
Use these only after identifying the topology and the port. Let $R_C'=R_C\parallel R_L$ and $R_E'=R_E\parallel R_L$ where relevant. These expressions assume $r_o=\infty$.
\begin{center}\small\begin{tabular}{p{2.7cm}p{4.4cm}p{7.3cm}}\toprule
Topology & Input resistance of transistor path & Loaded stage gain\\\midrule
CE, grounded emitter & $r_\pi$ & $-g_mR_C'$\\[7pt]
CE, unbypassed $R_E$ & $r_\pi+(\beta+1)R_E$ & $-\beta R_C'/[r_\pi+(\beta+1)R_E]$\\[7pt]
Emitter follower & $r_\pi+(\beta+1)R_E'$ & $(\beta+1)R_E'/[r_\pi+(\beta+1)R_E']$\\[7pt]
CB, grounded base & $r_\pi/(\beta+1)$ & $+g_mR_C'$\\[7pt]
CB, base grounded through $R_B$ & $(r_\pi+R_B)/(\beta+1)$ & $+\beta R_C'/(r_\pi+R_B)$\\\bottomrule
\end{tabular}\end{center}
For CE and follower input ports, put the bias resistance $R_B$ in parallel with the transistor path. For CB, any external emitter-to-ground resistor is in parallel with the transistor input resistance.

For CE and CB with $r_o=\infty$, $R_{\rm out}=R_C$. For a follower,
\[
R_{\rm out}=R_E\parallel\frac{r_\pi+(R_s\parallel R_B)}{\beta+1},
\]
where the external $R_L$ has been removed. If no base-bias resistor exists, use $R_s$ in place of $R_s\parallel R_B$.

\head{Worked warm-up: the three gains people mix up}
A grounded-emitter CE stage has $I_C=1$ mA, $\beta=100$, $R_C=2\K$, $R_L=2\K$, $R_B=10\K$, and $R_s=500\Oh$.
\[
g_m=40\mS,\quad r_\pi=2.5\K,\quad
R_{\rm in}=10\K\parallel2.5\K=2\K.
\]
\[
A_{v,\text{unloaded}}=-g_mR_C=-80,\quad
A_{v,\text{loaded}}=-g_m(R_C\parallel R_L)=-40.
\]
\[
\boxed{A_s=-40\frac{2000}{500+2000}=-32,\qquad R_{\rm out}=2\K.}
\]
\textbf{Interpretation:} the collector load halves the stage gain, and the source divider reduces it again. Output resistance is not $R_C\parallel R_L$ because $R_L$ is removed for that measurement.

\head{Two useful physical checks}
\textbf{Looking into the base:} emitter resistance is multiplied by $\beta+1$ because emitter current is $\beta+1$ times base current.\par
\textbf{Looking into the emitter:} base-side resistance is divided by $\beta+1$. This is why a follower has low output resistance and a CB stage has low input resistance.

\textbf{Source context:} \href{https://drive.google.com/drive/folders/1LX1CiD2Hl7SUBE9ZuGbItAZ45t7zIr8o}{Lectures 5--8} and \href{https://drive.google.com/drive/folders/1P1aq-hffhFVeBSUqwzJkBELg3kuHlFE1}{Sheets 6--7}. The source material guided the topics; the problems and solutions here were independently constructed.
\newpage
\head{Why the formulas work: replace the transistor}
The low-frequency hybrid-$\pi$ model below keeps the physical base, collector and emitter terminals. $v_\pi=v_b-v_e$ controls the current source; it is not always equal to the applied input voltage.
\begin{center}\begin{circuitikz}[scale=.9,transform shape]
\draw (0,3) node[above]{base $b$} to[short,o-] (1.5,3) to[R,l={$r_\pi$}] (1.5,0);
\draw (5,3) node[above]{collector $c$} to[cisource,l_={$g_mv_\pi$},o-] (5,0);
\draw (5,3) to[short,*-] (8,3) to[R,l={$r_o$}] (8,0);
\draw (1.5,0) -- (8,0);\draw (4,0) to[short,*-o] (4,-.7) node[below]{emitter $e$};
\node[left] at (1.5,2.6){$+$};\node[left] at (1.5,1.5){$v_\pi$};\node[left] at (1.5,.4){$-$};
\end{circuitikz}\end{center}
The diamond current source points from collector to emitter. If Early effect is neglected, remove the $r_o$ branch. Do not remove the controlled source when finding output resistance: its controlling voltage must be determined from the zeroed-source circuit.

\head{Derive one CE result instead of memorizing it}
Let an unbypassed resistor $R_E$ connect emitter to ground and assume $r_o=\infty$. Apply the input at the base. Then
\[
i_b=\frac{v_\pi}{r_\pi},\quad i_c=g_mv_\pi=\beta i_b,\quad i_e=(\beta+1)i_b.
\]
The input voltage must cover both the base--emitter voltage and the emitter-resistor drop:
\[
v_i=v_\pi+v_e=i_br_\pi+(\beta+1)i_bR_E.
\]
So the resistance of the transistor path, seen from the base, is
\[
R_{ib}=\frac{v_i}{i_b}=r_\pi+(\beta+1)R_E.
\]
The output at the collector is $v_o=-i_c(R_C\parallel R_L)$. Dividing by the input expression gives
\[
A_v=-\frac{\beta(R_C\parallel R_L)}{r_\pi+(\beta+1)R_E}.
\]
Only after this do you put $R_B$ in parallel with $R_{ib}$ to get the whole amplifier's input resistance. $R_s$ belongs outside that input port.

\head{A follower output-resistance measurement}
Remove the external load and set $v_s=0$. The base is connected to ground through $R_G=R_s\parallel R_B$, not necessarily through a short. A test current entering the emitter is shared between the internal $R_E$ and the transistor path. With $r_o=\infty$, that transistor path has resistance
\[
R_{e,\text{path}}=\frac{r_\pi+R_G}{\beta+1},\qquad
R_{\rm out}=R_E\parallel R_{e,\text{path}}.
\]
\textbf{Ask before every calculation:} Which terminal is the input? Which terminal is the output? Which terminal is at AC ground? Which resistor is the external load?

\newpage
\head{1. Grounded-emitter foundations}
\begin{center}\begin{circuitikz}[scale=.74,transform shape]
\draw (4,2.4) node[npn](q){$Q$};
\draw (-2,0) node[ground]{} to[sV,l={$v_s$}] (-2,2.4) to[R,l={$R_s$}] (0,2.4) -- (q.B);
\draw (1,2.4) node[above]{$v_i$} to[R,l={$R_B$},*-] (1,0) node[ground]{};
\draw (q.C) -- (4,4) to[R,l={$R_C$}] (4,6) node[ground,rotate=180]{};
\draw (4,4) to[short,*-] (7,4) node[above]{$v_o$} to[R,l={$R_L$}] (7,0) node[ground]{};
\draw (q.E) to[R,l={$R_E$}] (4,-.2) node[ground]{};\end{circuitikz}\par\small CE stage. $R_E=0$ means a wire; $R_B=\infty$ or $R_L=\infty$ means omit that branch.\end{center}\ex{1}{Start with a grounded emitter}$I_C=1$ mA, $\beta=100$, $R_C=2000\Oh$, $R_E=0$, $R_s=0$, $R_B=R_L=\infty$.
Find $g_m$, $r_\pi$, $R_{\rm in}$, $A_v$ and $R_{\rm out}$. State the phase relationship. For $v_i=1$ mV peak, find the output peak amplitude.\par\small\hyperref[sol:1]{Worked solution}\normalsize\work
\ex{2}{Halve the current, double the resistor}$I_C=0.5$ mA, $\beta=100$, $R_C=4000\Oh$, $R_E=0$, $R_s=0$, $R_B=R_L=\infty$.
Find $g_m$, $r_\pi$, $R_{\rm in}$, $A_v$ and $R_{\rm out}$. State the phase relationship. Compare your answers with Exercise 1: which quantities stay the same?\par\small\hyperref[sol:2]{Worked solution}\normalsize\work

\newpage
\head{2. Loading and Early effect}
\begin{center}\begin{circuitikz}[scale=.74,transform shape]
\draw (4,2.4) node[npn](q){$Q$};
\draw (-2,0) node[ground]{} to[sV,l={$v_s$}] (-2,2.4) to[R,l={$R_s$}] (0,2.4) -- (q.B);
\draw (1,2.4) node[above]{$v_i$} to[R,l={$R_B$},*-] (1,0) node[ground]{};
\draw (q.C) -- (4,4) to[R,l={$R_C$}] (4,6) node[ground,rotate=180]{};
\draw (4,4) to[short,*-] (7,4) node[above]{$v_o$} to[R,l={$R_L$}] (7,0) node[ground]{};
\draw (q.E) to[R,l={$R_E$}] (4,-.2) node[ground]{};\end{circuitikz}\par\small CE stage. $R_E=0$ means a wire; $R_B=\infty$ or $R_L=\infty$ means omit that branch.\end{center}\ex{3}{The source is not the base}$I_C=1$ mA, $\beta=100$, $R_C=3\K$, $R_L=6\K$, $R_B=15\K$, $R_s=1\K$, $R_E=0$. Find $g_m,r_\pi,R_{\rm in},A_v,A_s,R_{\rm out}$. For $v_s=2$ mV peak, find $v_i$ and the output peak amplitude.\par\small\hyperref[sol:3]{Worked solution}\normalsize\work
\ex{4}{Early effect with a grounded emitter}Repeat Exercise 3 with $V_A=20$ V. Use $r_o\simeq V_A/I_C$; retain the same bias current. Find $R_{\rm in},A_v,A_s,R_{\rm out}$ and explain why $r_o$ changes the collector-side results but not the input resistance here.\par\small\hyperref[sol:4]{Worked solution}\normalsize\work

\newpage
\head{3. Emitter degeneration}
\begin{center}\begin{circuitikz}[scale=.74,transform shape]
\draw (4,2.4) node[npn](q){$Q$};
\draw (-2,0) node[ground]{} to[sV,l={$v_s$}] (-2,2.4) to[R,l={$R_s$}] (0,2.4) -- (q.B);
\draw (1,2.4) node[above]{$v_i$} to[R,l={$R_B$},*-] (1,0) node[ground]{};
\draw (q.C) -- (4,4) to[R,l={$R_C$}] (4,6) node[ground,rotate=180]{};
\draw (4,4) to[short,*-] (7,4) node[above]{$v_o$} to[R,l={$R_L$}] (7,0) node[ground]{};
\draw (q.E) to[R,l={$R_E$}] (4,-.2) node[ground]{};\end{circuitikz}\par\small CE stage. $R_E=0$ means a wire; $R_B=\infty$ or $R_L=\infty$ means omit that branch.\end{center}\ex{5}{Add emitter degeneration}$I_C=1$ mA, $\beta=100$, $R_C=3\K$, $R_L=6\K$, $R_B=15\K$, $R_s=1\K$ and $R_E=100\Oh$. Derive $R_{ib}$ and $A_v$ from $v_i=i_br_\pi+i_eR_E$. Find $R_{\rm in},A_s,R_{\rm out}$. Compare $A_v$ with $-g_m(R_C\parallel R_L)/(1+g_mR_E)$.\par\small\hyperref[sol:5]{Worked solution}\normalsize\work
\ex{6}{Do not silently replace beta plus one}$I_C=1$ mA, $\beta=50$, $R_C=3\K$, $R_L=6\K$, $R_B=15\K$, $R_s=1\K$ and $R_E=470\Oh$. Derive $R_{ib}$ and $A_v$ from $v_i=i_br_\pi+i_eR_E$. Find $R_{\rm in},A_s,R_{\rm out}$. Compare $A_v$ with $-g_m(R_C\parallel R_L)/(1+g_mR_E)$.\par\small\hyperref[sol:6]{Worked solution}\normalsize\work

\newpage
\head{4. Bypass capacitors}
\begin{center}\begin{circuitikz}[scale=.74,transform shape]
\draw (4,2.4) node[npn](q){$Q$};
\draw (-2,0) node[ground]{} to[sV,l={$v_s$}] (-2,2.4) to[R,l={$R_s$}] (0,2.4) -- (q.B);
\draw (1,2.4) node[above]{$v_i$} to[R,l={$R_B$},*-] (1,0) node[ground]{};
\draw (q.C) -- (4,4) to[R,l={$R_C$}] (4,6) node[ground,rotate=180]{};
\draw (4,4) to[short,*-] (7,4) node[above]{$v_o$} to[R,l={$R_L$}] (7,0) node[ground]{};
\draw (q.E) -- (4,1.4) to[R,l_={$R_{E1}$}] (4,-.1) to[R,l_={$R_{E2}$}] (4,-1.7) node[ground]{};\draw (4,-.1) to[short,*-] (5.6,-.1) to[C,l={$C_E$}] (5.6,-1.7) node[ground]{};\end{circuitikz}\par\small CE stage. $R_E=0$ means a wire; $R_B=\infty$ or $R_L=\infty$ means omit that branch.\end{center}\ex{7}{Only part of the emitter resistor is bypassed}For the split-emitter figure, $I_C=1$ mA, $\beta=100$, $R_C=R_L=4\K$, $R_B=20\K$, $R_s=500\Oh$, $R_{E1}=100\Oh$ and $R_{E2}=900\Oh$. $C_E$ is a midband short. Find the emitter resistance used for AC and for DC. Then calculate $R_{\rm in},A_v,A_s,R_{\rm out}$.\par\small\hyperref[sol:7]{Worked solution}\normalsize\work
\ex{8}{What changes when the bypass capacitor is removed?}Use Exercise 7 but physically remove $C_E$. Recalculate $R_{\rm in},A_v,A_s,R_{\rm out}$. Does the DC operating point change? Explain why using $R_E=0$ would be wrong both before and after this removal.\par\small\hyperref[sol:8]{Worked solution}\normalsize\work

\newpage
\head{5. Emitter followers}
\begin{center}\begin{circuitikz}[scale=.75,transform shape]
\draw (4,2.8) node[npn](q){$Q$};\draw (q.C) -- (4,4.3) node[ground,rotate=180]{};
\draw (-2,0) node[ground]{} to[sV,l={$v_s$}] (-2,2.8) to[R,l={$R_s$}] (0,2.8) -- (q.B);
\draw (1,2.8) node[above]{$v_i$} to[R,l={$R_B$},*-] (1,0) node[ground]{};
\draw (q.E) -- (4,1.6) node[left]{$v_o$} to[R,l={$R_E$}] (4,-.7) node[ground]{};
\draw (4,1.6) to[short,*-] (6.5,1.6) to[R,l={$R_L$}] (6.5,-.7) node[ground]{};
\end{circuitikz}\par\small Emitter follower. The collector is at AC ground, not physically shorted to ground in DC.\end{center}\ex{9}{The basic emitter follower}$I_C=1$ mA, $\beta=100$, $R_E=R_L=1\K$, $R_B=\infty$, $R_s=0$. Find the transistor input resistance, loaded base-to-emitter gain, and output resistance with the external load removed. Explain why a gain below one does not make this stage useless.\par\small\hyperref[sol:9]{Worked solution}\normalsize\work
\ex{10}{Source resistance comes back during Rout}Use the same follower with $R_B=20\K$ and $R_s=3\K$. Find $R_{\rm in},A_v,A_s,R_{\rm out}$. A student claims $R_{\rm out}=1/g_m=25\Oh$. Explain what this estimate ignores.\par\small\hyperref[sol:10]{Worked solution}\normalsize\work

\newpage
\head{6. Common-base foundations}
\begin{center}\begin{circuitikz}[scale=.78,transform shape]
\draw (4,2.8) node[npn](q){$Q$};
\draw (q.B) to[R,l={$R_B$}] (0,2.8) node[ground]{};
\draw (q.C) -- (4,4) to[R,l={$R_C$}] (4,5.7) node[ground,rotate=180]{};
\draw (4,4) to[short,*-] (7,4) node[above]{$v_o$} to[R,l={$R_L$}] (7,.3) node[ground]{};
\draw (q.E) -- (4,1.5) to[R,l={$R_E$}] (4,-.7) node[ground]{};
\draw (-2,-.7) node[ground]{} to[sV,l={$v_s$}] (-2,1.5) to[R,l={$R_s$}] (1,1.5) -- (4,1.5) node[above right]{$v_i$};
\end{circuitikz}\par\small Common-base stage. $R_B=0$ shorts the base to AC ground; $R_E=\infty$ omits the emitter resistor.\end{center}\ex{11}{Common base: input at the emitter}$I_C=1$ mA, $\beta=100$, $R_C=2\K$, $R_L=\infty$, $R_E=\infty$, $R_B=0$, $R_s=0$. Derive the emitter input resistance and voltage gain. Explain why a positive emitter-voltage change raises the collector voltage.\par\small\hyperref[sol:11]{Worked solution}\normalsize\work
\ex{12}{A high stage gain can hide a small overall gain}$I_C=1$ mA, $\beta=100$, $R_C=R_L=2\K$, $R_E=1\K$, $R_B=0$, $R_s=100\Oh$. Find $R_{\rm in},A_v,A_s,R_{\rm out}$. Explain why $A_s$ is much smaller than $A_v$.\par\small\hyperref[sol:12]{Worked solution}\normalsize\work

\newpage
\head{7. Finite base resistance}
\begin{center}\begin{circuitikz}[scale=.78,transform shape]
\draw (4,2.8) node[npn](q){$Q$};
\draw (q.B) to[R,l={$R_B$}] (0,2.8) node[ground]{};
\draw (q.C) -- (4,4) to[R,l={$R_C$}] (4,5.7) node[ground,rotate=180]{};
\draw (4,4) to[short,*-] (7,4) node[above]{$v_o$} to[R,l={$R_L$}] (7,.3) node[ground]{};
\draw (q.E) -- (4,1.5) to[R,l={$R_E$}] (4,-.7) node[ground]{};
\draw (-2,-.7) node[ground]{} to[sV,l={$v_s$}] (-2,1.5) to[R,l={$R_s$}] (1,1.5) -- (4,1.5) node[above right]{$v_i$};
\end{circuitikz}\par\small Common-base stage. $R_B=0$ shorts the base to AC ground; $R_E=\infty$ omits the emitter resistor.\end{center}\ex{13}{The base is grounded through a resistor}$I_C=1$ mA, $\beta=100$, $R_C=R_L=2\K$, $R_E=\infty$, $R_B=2.5\K$, $R_s=100\Oh$. Derive $v_\pi/v_i$, then find $R_{\rm in},A_v,A_s,R_{\rm out}$. Why is it wrong to set $v_b=0$?\par\small\hyperref[sol:13]{Worked solution}\normalsize\work
\ex{14}{Compare two base terminations}Use $I_C=1$ mA, $\beta=100$, $R_C=R_L=2\K$, $R_E=1\K$, $R_s=100\Oh$. Calculate $R_{\rm in},A_v,A_s$ for (i) $R_B=0$ and (ii) $R_B=10\K$. Decide whether the increased emitter input resistance necessarily improves the overall voltage gain.\par\small\hyperref[sol:14]{Worked solution}\normalsize\work

\newpage
\head{8. CE--follower cascades}
\begin{center}\begin{circuitikz}[scale=.67,transform shape]
\draw (1,2.8) node[npn](q1){$Q_1$};\draw (q1.E) -- (1,1) node[ground]{};
\draw (-4,.6) node[ground]{} to[sV,l={$v_s$}] (-4,2.8) to[R,l={$R_s$}] (-2,2.8) -- (q1.B);
\draw (-1,2.8) node[above]{$v_i$} to[R,l={$R_{B1}$},*-] (-1,.6) node[ground]{};
\draw (q1.C) -- (1,4.2) to[R,l={$R_{C1}$}] (1,6) node[ground,rotate=180]{};
\draw (7,2.8) node[npn](q2){$Q_2$};
\draw (1,4.2) to[short,*-] (4,4.2) node[above]{$v_x$} -- (4,2.8) -- (q2.B);
\draw (4,2.8) to[R,l={$R_{B2}$},*-] (4,.6) node[ground]{};
\draw (q2.C) -- (7,5.5) node[ground,rotate=180]{};\draw (q2.E) -- (7,1.5) to[R,l={$R_{E2}$}] (7,-.6) node[ground]{};\draw (7,1.5) to[short,*-] (9.5,1.5) node[above]{$v_o$} to[R,l={$R_L$}] (9.5,-.6) node[ground]{};\end{circuitikz}\par\small Midband cascade: interstage coupling capacitor is replaced by a short. Biases are set independently.\end{center}\ex{15}{CE driving a follower}For the CE--follower cascade, $I_{C1}=1$ mA, $I_{C2}=2$ mA, $\beta_1=\beta_2=100$, $R_{B1}=20\K$, $R_{B2}=50\K$, $R_{C1}=3\K$, $R_{E2}=1\K$, $R_L=500\Oh$, and $R_s=1\K$. Find $R_{\rm in1},R_{\rm in2},A_1=v_x/v_i,A_2=v_o/v_x,A_v,A_s$, and the complete cascade's output resistance.\par\small\hyperref[sol:15]{Worked solution}\normalsize\work
\ex{16}{Change the final load and follow the effect backward}Repeat Exercise 15 with $R_L=50\Oh$. Explain why the first-stage gain also changes, even though $R_{C1}$ and $I_{C1}$ did not change. Does the measured output resistance change when the external load value changes?\par\small\hyperref[sol:16]{Worked solution}\normalsize\work

\newpage
\head{9. CE--CE cascades}
\begin{center}\begin{circuitikz}[scale=.67,transform shape]
\draw (1,2.8) node[npn](q1){$Q_1$};\draw (q1.E) -- (1,1) node[ground]{};
\draw (-4,.6) node[ground]{} to[sV,l={$v_s$}] (-4,2.8) to[R,l={$R_s$}] (-2,2.8) -- (q1.B);
\draw (-1,2.8) node[above]{$v_i$} to[R,l={$R_{B1}$},*-] (-1,.6) node[ground]{};
\draw (q1.C) -- (1,4.2) to[R,l={$R_{C1}$}] (1,6) node[ground,rotate=180]{};
\draw (7,2.8) node[npn](q2){$Q_2$};
\draw (1,4.2) to[short,*-] (4,4.2) node[above]{$v_x$} -- (4,2.8) -- (q2.B);
\draw (4,2.8) to[R,l={$R_{B2}$},*-] (4,.6) node[ground]{};
\draw (q2.E) -- (7,1) node[ground]{};\draw (q2.C) -- (7,4.2) to[R,l={$R_{C2}$}] (7,6) node[ground,rotate=180]{};\draw (7,4.2) to[short,*-] (9.5,4.2) node[above]{$v_o$} to[R,l={$R_L$}] (9.5,1) node[ground]{};\end{circuitikz}\par\small Midband cascade: interstage coupling capacitor is replaced by a short. Biases are set independently.\end{center}\ex{17}{Two CE stages: do not multiply unloaded gains}In the CE--CE cascade, $I_{C1}=I_{C2}=1$ mA, $\beta_1=\beta_2=100$, $R_{B1}=20\K$, $R_{B2}=10\K$, $R_{C1}=3\K$, $R_{C2}=R_L=2\K$, $R_s=500\Oh$. Find both input resistances, the loaded gain of each stage, $A_v$, $A_s$ and output resistance. Compare $A_v$ with the product of both stages' unloaded gains.\par\small\hyperref[sol:17]{Worked solution}\normalsize\work
\ex{18}{An extra interstage resistor}Use Exercise 17 but insert $R_x=1\K$ in the interstage wire, between the first collector ($v_x$) and the second base ($v_{b2}$). Find $v_x/v_i$, $v_{b2}/v_x$, $v_o/v_{b2}$, $A_v$ and $A_s$. Is the first-stage gain alone sufficient to judge the overall improvement?\par\small\hyperref[sol:18]{Worked solution}\normalsize\work

\newpage
\head{10. DC bias to AC performance}
\begin{center}\begin{circuitikz}[scale=.75,transform shape]
\draw (4,3) node[npn](q){$Q$};
\draw (1,6) -- (4,6) node[above]{$8\,\mathrm V$} to[R,l={$3\K$}] (q.C);
\draw (1,6) to[R,l_={$60\K$}] (1,3) -- (q.B);
\draw (1,3) to[R,l_={$20\K$},*-] (1,0) node[ground]{};
\draw (-2,3) node[left]{$v_i$} to[C,l={$C_{\rm in}$}] (1,3);
\draw (q.E) -- (4,1.9) to[R,l_={$1\K$}] (4,0) node[ground]{};
\draw (4,1.9) to[short,*-] (5.8,1.9) to[C,l={$C_E$}] (5.8,0) node[ground]{};
\draw (q.C) to[short,*-] (6.5,3.77) to[C,l={$C_{\rm out}$}] (8,3.77) node[above]{$v_o$} to[R,l={$6\K$}] (8,0) node[ground]{};
\end{circuitikz}\par\small DC circuit for Exercises 19--20. All capacitors are large for midband AC.\end{center}\ex{19}{From the DC circuit to the AC circuit}Use the DC figure with $\beta=100$, $V_{BE}=0.70$ V, $V_A=\infty$. The AC source driving $C_{\rm in}$ is ideal. Calculate $I_C,V_B,V_E,V_C,V_{CE}$ and check forward-active operation. Draw the midband model and calculate $g_m,r_\pi,R_{\rm in},A_v,R_{\rm out}$ with $C_E$ connected.\par\small\hyperref[sol:19]{Worked solution}\normalsize\work
\ex{20}{Remove the bypass and add a realistic source}In Exercise 19, remove $C_E$ and drive $C_{\rm in}$ from a source with $R_s=1\K$. Calculate the new $R_{\rm in},A_v,A_s,R_{\rm out}$. Explain why neither removing $C_E$ nor changing the AC source resistance changes the DC bias in this circuit.\par\small\hyperref[sol:20]{Worked solution}\normalsize\work

\newpage
\head{11. Finite-output-resistance challenges}
Use the CE circuit from page \pageref{ex:5} for Exercise 21 and the follower circuit from page \pageref{ex:9} for Exercise 22. Those figures retain their labeled topology; the values below replace earlier values.\par
\textbf{Focus:} a finite $r_o$ connects collector to emitter. It is a resistor to ground only when the opposite terminal is actually at AC ground.\vspace{12pt}\ex{21}{Challenge: Rout with degeneration and finite ro}Use the CE figure with $I_C=1$ mA, $\beta=100$, $V_A=50$ V, $R_C=4\K$, $R_E=200\Oh$, $R_B=\infty$, and $R_s=0$. Remove $R_L$. Find $r_o$ and derive the resistance looking into the transistor collector using a test voltage, then include $R_C$. Compare with the incorrect shortcut $R_C\parallel r_o$. Do not assume the emitter is at AC ground.\par\small\hyperref[sol:21]{Worked solution}\normalsize\work
\ex{22}{Challenge: finite ro in a follower}Use the emitter-follower figure with $I_C=1$ mA, $\beta=100$, $r_o=2\K$ (given directly), $R_E=R_L=1\K$, $R_B=20\K$, $R_s=3\K$. Find $R_{\rm in},A_v,A_s,R_{\rm out}$. Explain where $r_o$ appears in the AC circuit.\par\small\hyperref[sol:22]{Worked solution}\normalsize\work

\newpage
\head{12. Design and error detection}
\begin{center}\begin{circuitikz}[scale=.74,transform shape]
\draw (4,2.4) node[npn](q){$Q$};
\draw (-2,0) node[ground]{} to[sV,l={$v_s$}] (-2,2.4) to[R,l={$R_s$}] (0,2.4) -- (q.B);
\draw (1,2.4) node[above]{$v_i$} to[R,l={$R_B$},*-] (1,0) node[ground]{};
\draw (q.C) -- (4,4) to[R,l={$R_C$}] (4,6) node[ground,rotate=180]{};
\draw (4,4) to[short,*-] (7,4) node[above]{$v_o$} to[R,l={$R_L$}] (7,0) node[ground]{};
\draw (q.E) to[R,l={$R_E$}] (4,-.2) node[ground]{};\end{circuitikz}\par\small CE stage. $R_E=0$ means a wire; $R_B=\infty$ or $R_L=\infty$ means omit that branch.\end{center}\ex{23}{Design for a specified loaded gain}Use the CE figure with $I_C=1$ mA, $\beta=100$, $R_C=R_L=4\K$, $R_B=50\K$, $R_s=1\K$, and $r_o=\infty$. Select the unbypassed $R_E$ so the \emph{base-to-output loaded gain} is exactly $-10$. Then find $R_{\rm in},A_s,R_{\rm out}$. Explain why designing $A_v=-10$ does not make $A_s=-10$. The bias circuit is adjusted to preserve $I_C$.\par\small\hyperref[sol:23]{Worked solution}\normalsize\work
\ex{24}{Diagnose a complete wrong solution}Use $I_C=1$ mA, $\beta=100$, $R_C=3\K$, $R_L=6\K$, $R_B=15\K$, $R_s=1\K$, $R_E=100\Oh$, $r_o=\infty$. A student writes:
\[R_{\rm in}=R_s+r_\pi,\quad A_v=-g_mR_C,\quad A_s=A_v,\quad R_{\rm out}=R_C\parallel R_L.\]
Identify every conceptual mistake, derive the corrected quantities, and calculate the output amplitude for $v_s=1$ mV peak.\par\small\hyperref[sol:24]{Worked solution}\normalsize\work
\newpage\label{solutions}\head{Worked solutions}\textit{Use these to check the circuit reasoning, not only the final numbers. Rounding to three significant figures is sufficient.}\par\sol{1}{Start with a grounded emitter}First $g_m=40\mS$ and $r_\pi=2500\Oh$. The transistor input path is
\[R_{ib}=r_\pi+(\beta+1)R_E=2500\Oh,\quad R_{\rm in}=R_B\parallel R_{ib}=2500\Oh.\]
The effective collector load is $R_C\parallel R_L=2000\Oh$. Hence
\[A_v=-\frac{\beta(2000)}{2500+(\beta+1)(0)}=-80,\qquad
A_s=A_v\frac{2500}{0+2500}=-80.\]
With the load removed and the source zeroed, $R_{\rm out}=2000\Oh$. The output is inverted and its peak amplitude is $80$ mV.\par\vfill
\sol{2}{Halve the current, double the resistor}First $g_m=20\mS$ and $r_\pi=5000\Oh$. The transistor input path is
\[R_{ib}=r_\pi+(\beta+1)R_E=5000\Oh,\quad R_{\rm in}=R_B\parallel R_{ib}=5000\Oh.\]
The effective collector load is $R_C\parallel R_L=4000\Oh$. Hence
\[A_v=-\frac{\beta(4000)}{5000+(\beta+1)(0)}=-80,\qquad
A_s=A_v\frac{5000}{0+5000}=-80.\]
With the load removed and the source zeroed, $R_{\rm out}=4000\Oh$. The gain stays $-80$; $r_pi$, input resistance and output resistance double relative to Exercise 1. The output is inverted.\par\vfill

\newpage
\head{Worked solutions: 3--4}
\sol{3}{The source is not the base}First $g_m=40\mS$ and $r_\pi=2500\Oh$. The transistor input path is
\[R_{ib}=r_\pi+(\beta+1)R_E=2500\Oh,\quad R_{\rm in}=R_B\parallel R_{ib}=2143\Oh.\]
The effective collector load is $R_C\parallel R_L=2000\Oh$. Hence
\[A_v=-\frac{\beta(2000)}{2500+(\beta+1)(0)}=-80,\qquad
A_s=A_v\frac{2143}{1000+2143}=-54.55.\]
With the load removed and the source zeroed, $R_{\rm out}=3000\Oh$. The input fraction is $0.6818$, so the base signal is $1.364$ mV peak and the output amplitude is $109.1$ mV peak. $R_L$ affects gain but is removed for $R_{\rm out}$.\par\vfill
\sol{4}{Early effect with a grounded emitter}First $g_m=40\mS$ and $r_\pi=2500\Oh$. The transistor input path is
\[R_{ib}=r_\pi+(\beta+1)R_E=2500\Oh,\quad R_{\rm in}=R_B\parallel R_{ib}=2143\Oh.\]
The effective collector load is $R_C\parallel R_L\parallel r_o=1818\Oh$. Hence
\[A_v=-\frac{\beta(1818)}{2500+(\beta+1)(0)}=-72.73,\qquad
A_s=A_v\frac{2143}{1000+2143}=-49.59.\]
With the load removed and the source zeroed, $R_{\rm out}=2609\Oh$. With emitter at AC ground, $r_o$ is another collector-to-ground resistor. It introduces no base-to-collector feedback in this low-frequency hybrid-$\pi$ model, so $R_{\rm in}$ is unchanged.\par\vfill

\newpage
\head{Worked solutions: 5--6}
\sol{5}{Add emitter degeneration}First $g_m=40\mS$ and $r_\pi=2500\Oh$. The transistor input path is
\[R_{ib}=r_\pi+(\beta+1)R_E=12600\Oh,\quad R_{\rm in}=R_B\parallel R_{ib}=6848\Oh.\]
The effective collector load is $R_C\parallel R_L=2000\Oh$. Hence
\[A_v=-\frac{\beta(2000)}{2500+(\beta+1)(100)}=-15.87,\qquad
A_s=A_v\frac{6848}{1000+6848}=-13.85.\]
With the load removed and the source zeroed, $R_{\rm out}=3000\Oh$. The exact relation is $i_e=(\beta+1)i_b$, giving the expressions above. The large-$\beta$ gain approximation is $-16$. Degeneration raises input resistance and lowers gain magnitude; with $r_o=\infty$, it does not change $R_{\rm out}=R_C$.\par\vfill
\sol{6}{Do not silently replace beta plus one}First $g_m=40\mS$ and $r_\pi=1250\Oh$. The transistor input path is
\[R_{ib}=r_\pi+(\beta+1)R_E=25220\Oh,\quad R_{\rm in}=R_B\parallel R_{ib}=9406\Oh.\]
The effective collector load is $R_C\parallel R_L=2000\Oh$. Hence
\[A_v=-\frac{\beta(2000)}{1250+(\beta+1)(470)}=-3.965,\qquad
A_s=A_v\frac{9406}{1000+9406}=-3.584.\]
With the load removed and the source zeroed, $R_{\rm out}=3000\Oh$. The exact relation is $i_e=(\beta+1)i_b$, giving the expressions above. The large-$\beta$ gain approximation is $-4.04$. Degeneration raises input resistance and lowers gain magnitude; with $r_o=\infty$, it does not change $R_{\rm out}=R_C$.\par\vfill

\newpage
\head{Worked solutions: 7--8}
\sol{7}{Only part of the emitter resistor is bypassed}For DC, $C_E$ is open and the emitter resistance is $100+900=1000\Oh$. For midband AC, it shorts only $R_{E2}$, leaving $R_E=100\Oh$. First $g_m=40\mS$ and $r_\pi=2500\Oh$. The transistor input path is
\[R_{ib}=r_\pi+(\beta+1)R_E=12600\Oh,\quad R_{\rm in}=R_B\parallel R_{ib}=7730\Oh.\]
The effective collector load is $R_C\parallel R_L=2000\Oh$. Hence
\[A_v=-\frac{\beta(2000)}{2500+(\beta+1)(100)}=-15.87,\qquad
A_s=A_v\frac{7730}{500+7730}=-14.91.\]
With the load removed and the source zeroed, $R_{\rm out}=4000\Oh$.\par\vfill
\sol{8}{What changes when the bypass capacitor is removed?}First $g_m=40\mS$ and $r_\pi=2500\Oh$. The transistor input path is
\[R_{ib}=r_\pi+(\beta+1)R_E=103500\Oh,\quad R_{\rm in}=R_B\parallel R_{ib}=16760\Oh.\]
The effective collector load is $R_C\parallel R_L=2000\Oh$. Hence
\[A_v=-\frac{\beta(2000)}{2500+(\beta+1)(1000)}=-1.932,\qquad
A_s=A_v\frac{16760}{500+16760}=-1.876.\]
With the load removed and the source zeroed, $R_{\rm out}=4000\Oh$. The DC bias is unchanged because the capacitor was already open for DC. Originally $R_{E1}$ remained unbypassed; after removal both resistors remain in AC. A capacitor bypasses only the component(s) across which it is connected.\par\vfill

\newpage
\head{Worked solutions: 9--10}
\sol{9}{The basic emitter follower}$g_m=40\mS$, $r_\pi=2500\Oh$. The loaded emitter resistance is $R_E'=500\Oh$.
\[R_{ib}=r_\pi+(\beta+1)R_E'=53000\Oh,\quad R_{\rm in}=R_B\parallel R_{ib}=53000\Oh.\]
\[A_v=\frac{(\beta+1)R_E'}{r_\pi+(\beta+1)R_E'}=0.9528,\quad A_s=A_v\frac{R_{\rm in}}{R_s+R_{\rm in}}=0.9528.\]
For output resistance, remove $R_L$ and zero $v_s$. The base now sees $R_s\parallel R_B=0\Oh$:
\[R_{\rm out}=R_E\parallel\frac{r_\pi+(R_s\parallel R_B)}{\beta+1}=24.15\Oh.\]
The gain is positive and less than unity; the follower supplies impedance transformation.\par\vfill
\sol{10}{Source resistance comes back during Rout}$g_m=40\mS$, $r_\pi=2500\Oh$. The loaded emitter resistance is $R_E'=500\Oh$.
\[R_{ib}=r_\pi+(\beta+1)R_E'=53000\Oh,\quad R_{\rm in}=R_B\parallel R_{ib}=14520\Oh.\]
\[A_v=\frac{(\beta+1)R_E'}{r_\pi+(\beta+1)R_E'}=0.9528,\quad A_s=A_v\frac{R_{\rm in}}{R_s+R_{\rm in}}=0.7897.\]
For output resistance, remove $R_L$ and zero $v_s$. The base now sees $R_s\parallel R_B=2609\Oh$:
\[R_{\rm out}=R_E\parallel\frac{r_\pi+(R_s\parallel R_B)}{\beta+1}=48.15\Oh.\]
The gain is positive and less than unity; the follower supplies impedance transformation. The $1/g_m$ shortcut ignores base-side source/bias resistance and the parallel internal emitter resistor. It is only an approximation for a strongly driven base and sufficiently large $R_E$.\par\vfill

\newpage
\head{Worked solutions: 11--12}
\sol{11}{Common base: input at the emitter}$g_m=40\mS$, $r_\pi=2500\Oh$. Looking into the transistor emitter,
\[R_{ie}=\frac{r_\pi+R_B}{\beta+1}=24.75\Oh,\quad R_{\rm in}=R_E\parallel R_{ie}=24.75\Oh.\]
Since $v_\pi=-v_i r_\pi/(r_\pi+R_B)$,
\[A_v=+\frac{\beta(R_C\parallel R_L)}{r_\pi+R_B}=80,\quad A_s=A_v\frac{R_{\rm in}}{R_s+R_{\rm in}}=80.\]
$R_{\rm out}=R_C=2000\Oh$ because $r_o=\infty$. The stage is noninverting. Raising $v_e$ with fixed base voltage reduces $v_{BE}$ and collector current. The collector resistor then drops less voltage, so $v_o$ rises.\par\vfill
\sol{12}{A high stage gain can hide a small overall gain}$g_m=40\mS$, $r_\pi=2500\Oh$. Looking into the transistor emitter,
\[R_{ie}=\frac{r_\pi+R_B}{\beta+1}=24.75\Oh,\quad R_{\rm in}=R_E\parallel R_{ie}=24.15\Oh.\]
Since $v_\pi=-v_i r_\pi/(r_\pi+R_B)$,
\[A_v=+\frac{\beta(R_C\parallel R_L)}{r_\pi+R_B}=40,\quad A_s=A_v\frac{R_{\rm in}}{R_s+R_{\rm in}}=7.782.\]
$R_{\rm out}=R_C=2000\Oh$ because $r_o=\infty$. The stage is noninverting. The source resistance is large compared with the roughly $25\Oh$ input resistance, so most source voltage is lost before reaching the emitter.\par\vfill

\newpage
\head{Worked solutions: 13--14}
\sol{13}{The base is grounded through a resistor}$g_m=40\mS$, $r_\pi=2500\Oh$. Looking into the transistor emitter,
\[R_{ie}=\frac{r_\pi+R_B}{\beta+1}=49.5\Oh,\quad R_{\rm in}=R_E\parallel R_{ie}=49.5\Oh.\]
Since $v_\pi=-v_i r_\pi/(r_\pi+R_B)$,
\[A_v=+\frac{\beta(R_C\parallel R_L)}{r_\pi+R_B}=20,\quad A_s=A_v\frac{R_{\rm in}}{R_s+R_{\rm in}}=6.623.\]
$R_{\rm out}=R_C=2000\Oh$ because $r_o=\infty$. The stage is noninverting. $v_\pi/v_i=-1/2$. The finite base resistor carries small-signal base current and therefore develops a nonzero voltage. Only an ideal AC ground fixes $v_b=0$.\par\vfill
\sol{14}{Compare two base terminations}\textbf{(i) Grounded base.} $g_m=40\mS$, $r_\pi=2500\Oh$. Looking into the transistor emitter,
\[R_{ie}=\frac{r_\pi+R_B}{\beta+1}=24.75\Oh,\quad R_{\rm in}=R_E\parallel R_{ie}=24.15\Oh.\]
Since $v_\pi=-v_i r_\pi/(r_\pi+R_B)$,
\[A_v=+\frac{\beta(R_C\parallel R_L)}{r_\pi+R_B}=40,\quad A_s=A_v\frac{R_{\rm in}}{R_s+R_{\rm in}}=7.782.\]
$R_{\rm out}=R_C=2000\Oh$ because $r_o=\infty$. The stage is noninverting.\par\textbf{(ii) $R_B=10\K$.} $g_m=40\mS$, $r_\pi=2500\Oh$. Looking into the transistor emitter,
\[R_{ie}=\frac{r_\pi+R_B}{\beta+1}=123.8\Oh,\quad R_{\rm in}=R_E\parallel R_{ie}=110.1\Oh.\]
Since $v_\pi=-v_i r_\pi/(r_\pi+R_B)$,
\[A_v=+\frac{\beta(R_C\parallel R_L)}{r_\pi+R_B}=8,\quad A_s=A_v\frac{R_{\rm in}}{R_s+R_{\rm in}}=4.193.\]
$R_{\rm out}=R_C=2000\Oh$ because $r_o=\infty$. The stage is noninverting. The source divider improves, but the base moves with the emitter, reducing $|v_\pi/v_i|$. Overall gain falls; input resistance alone cannot predict gain.\par\vfill

\newpage
\head{Worked solutions: 15--16}
\sol{15}{CE driving a follower}Stage 1 has $g_{m1}=40\mS$, $r_{\pi1}=2.5\K$, and $R_{\rm in1}=20\K\parallel2.5\K=2222\Oh$.
For stage 2, $g_{m2}=80\mS$, $r_{\pi2}=1.25\K$, and $R_E'=1\K\parallel R_L=333.3\Oh$.
\[R_{\rm in2}=50\K\parallel[1250+101R_E']=20560\Oh.\]
Stage 2 loads stage 1. Therefore
\[A_1=\frac{v_x}{v_i}=-g_{m1}(3\K\parallel R_{\rm in2})=-104.7,\quad
A_2=\frac{v_o}{v_x}=\frac{101R_E'}{1250+101R_E'}=0.9642.\]
\[A_v=A_1A_2=-101,\qquad A_s=A_v\frac{R_{\rm in1}}{1000+R_{\rm in1}}=-69.63.\]
For the final output resistance, remove $R_L$. With $v_s=0$, the follower base sees $R_{C1}\parallel R_{B2}=3\K\parallel50\K$. Thus
\[R_{\rm out}=1\K\parallel\frac{1250+(3\K\parallel50\K)}{101}=38.83\Oh.\]
Do not use the loaded $R_{\rm in2}$ as the source resistance at its own base.\par\vfill
\sol{16}{Change the final load and follow the effect backward}Stage 1 has $g_{m1}=40\mS$, $r_{\pi1}=2.5\K$, and $R_{\rm in1}=20\K\parallel2.5\K=2222\Oh$.
For stage 2, $g_{m2}=80\mS$, $r_{\pi2}=1.25\K$, and $R_E'=1\K\parallel R_L=47.62\Oh$.
\[R_{\rm in2}=50\K\parallel[1250+101R_E']=5405\Oh.\]
Stage 2 loads stage 1. Therefore
\[A_1=\frac{v_x}{v_i}=-g_{m1}(3\K\parallel R_{\rm in2})=-77.17,\quad
A_2=\frac{v_o}{v_x}=\frac{101R_E'}{1250+101R_E'}=0.7937.\]
\[A_v=A_1A_2=-61.25,\qquad A_s=A_v\frac{R_{\rm in1}}{1000+R_{\rm in1}}=-42.24.\]
For the final output resistance, remove $R_L$. With $v_s=0$, the follower base sees $R_{C1}\parallel R_{B2}=3\K\parallel50\K$. Thus
\[R_{\rm out}=1\K\parallel\frac{1250+(3\K\parallel50\K)}{101}=38.83\Oh.\]
Do not use the loaded $R_{\rm in2}$ as the source resistance at its own base. Lower $R_L$ lowers the follower's input resistance, which loads the CE collector more heavily. Output resistance is unchanged here because the external load is removed during that measurement and the bias currents are held fixed.\par\vfill

\newpage
\head{Worked solutions: 17--18}
\sol{17}{Two CE stages: do not multiply unloaded gains}Both transistors have $g_m=40\mS$ and $r_\pi=2.5\K$.
\[R_{\rm in1}=20\K\parallel2.5\K=2222\Oh,\qquad R_{\rm in2}=10\K\parallel2.5\K=2\K.\]
The first collector sees $R_{C1}\parallel(R_x+R_{\rm in2})$, with $R_x=0\Oh$:
\[\frac{v_x}{v_i}=-.04[3000\parallel(0+2000)]=-48.\]
\[\frac{v_{b2}}{v_x}=\frac{2000}{0+2000}=1,\quad
\frac{v_o}{v_{b2}}=-.04(2000\parallel2000)=-40.\]
\[A_v=1920,\qquad A_s=A_v\frac{2222}{500+2222}=1567.\]
The two inversions produce a positive total gain. $R_{\rm out}=R_{C2}=2\K$ after removal of $R_L$. The unloaded-gain product would be $(-120)(-80)=9600$, while the actual base-input-to-final-output gain is $1920$. Both the interstage load and final load must be included.\par\vfill
\sol{18}{An extra interstage resistor}Both transistors have $g_m=40\mS$ and $r_\pi=2.5\K$.
\[R_{\rm in1}=20\K\parallel2.5\K=2222\Oh,\qquad R_{\rm in2}=10\K\parallel2.5\K=2\K.\]
The first collector sees $R_{C1}\parallel(R_x+R_{\rm in2})$, with $R_x=1000\Oh$:
\[\frac{v_x}{v_i}=-.04[3000\parallel(1000+2000)]=-60.\]
\[\frac{v_{b2}}{v_x}=\frac{2000}{1000+2000}=0.6667,\quad
\frac{v_o}{v_{b2}}=-.04(2000\parallel2000)=-40.\]
\[A_v=1600,\qquad A_s=A_v\frac{2222}{500+2222}=1306.\]
The two inversions produce a positive total gain. $R_{\rm out}=R_{C2}=2\K$ after removal of $R_L$. The first collector is less heavily loaded, so its local gain magnitude rises. But $R_x$ introduces a second divider; the overall gain falls. Trace the voltage through every intermediate node.\par\vfill

\newpage
\head{Worked solutions: 19--20}
\sol{19}{From the DC circuit to the AC circuit}The DC Thevenin source is $V_{th}=8(20/80)=2$ V and $R_{th}=60\K\parallel20\K=15\K$.
\[I_B=\frac{2-.7}{15000+101(1000)}=11.21\,\mu\mathrm A,\quad I_C=1.121\,\mathrm{mA}.\]
\[V_E=1.132\,\mathrm V,\quad V_B=1.832\,\mathrm V,\quad V_C=4.638\,\mathrm V.\]
$V_{CE}=3.506$ V and $V_C>V_B$, so forward-active operation is consistent. For midband, the emitter is grounded by $C_E$, the supply is AC ground, and the bias resistors become $15\K$ from base to ground. Draw $r_\pi$ from base to emitter and a controlled current $g_mv_\pi$ from collector to emitter. First $g_m=44.83\mS$ and $r_\pi=2231\Oh$. The transistor input path is
\[R_{ib}=r_\pi+(\beta+1)R_E=2231\Oh,\quad R_{\rm in}=R_B\parallel R_{ib}=1942\Oh.\]
The effective collector load is $R_C\parallel R_L=2000\Oh$. Hence
\[A_v=-\frac{\beta(2000)}{2231+(\beta+1)(0)}=-89.66,\qquad
A_s=A_v\frac{1942}{0+1942}=-89.66.\]
With the load removed and the source zeroed, $R_{\rm out}=3000\Oh$.\par\vfill
\sol{20}{Remove the bypass and add a realistic source}First $g_m=44.83\mS$ and $r_\pi=2231\Oh$. The transistor input path is
\[R_{ib}=r_\pi+(\beta+1)R_E=103200\Oh,\quad R_{\rm in}=R_B\parallel R_{ib}=13100\Oh.\]
The effective collector load is $R_C\parallel R_L=2000\Oh$. Hence
\[A_v=-\frac{\beta(2000)}{2231+(\beta+1)(1000)}=-1.937,\qquad
A_s=A_v\frac{13100}{1000+13100}=-1.8.\]
With the load removed and the source zeroed, $R_{\rm out}=3000\Oh$. The capacitor was already open in DC, and $C_{\rm in}$ isolates the source resistance from the bias network. Thus the bias current and small-signal parameters remain those of Exercise 19.\par\vfill

\newpage
\head{Worked solutions: 21--22}
\sol{21}{Challenge: Rout with degeneration and finite ro}$g_m=.04$ S, $r_\pi=2500\Oh$, $r_o=50000\Oh$. The zeroed ideal source grounds the base, so let $E=R_E\parallel r_\pi=185.2\Oh$. Apply $v_x$ to the collector. Since $v_\pi=-v_e$, emitter KCL gives
\[\frac{v_e}{E}+\frac{v_e-v_x}{r_o}+g_mv_e=0,
\quad v_e=\frac{v_x}{1+r_o/E+g_mr_o}.\]
The test current into the transistor is $i_x=(v_x-v_e)/r_o-g_mv_e$. Substitution gives
\[R_{\rm out,Q}=r_o+(1+g_mr_o)(R_E\parallel r_\pi)=420600\Oh.\]
\[R_{\rm out}=R_C\parallel R_{\rm out,Q}=3962\Oh.\]
The shortcut $4\K\parallel50\K=3704\Oh$ ignores the emitter motion and the feedback created by degeneration. The idealized $r_o=\infty$ answer would be $4\K$.\par\vfill
\sol{22}{Challenge: finite ro in a follower}$g_m=40\mS$, $r_\pi=2500\Oh$. The loaded emitter resistance is $R_E'=400\Oh$.
\[R_{ib}=r_\pi+(\beta+1)R_E'=42900\Oh,\quad R_{\rm in}=R_B\parallel R_{ib}=13640\Oh.\]
\[A_v=\frac{(\beta+1)R_E'}{r_\pi+(\beta+1)R_E'}=0.9417,\quad A_s=A_v\frac{R_{\rm in}}{R_s+R_{\rm in}}=0.7719.\]
For output resistance, remove $R_L$ and zero $v_s$. The base now sees $R_s\parallel R_B=2609\Oh$:
\[R_{\rm out}=R_E\parallel r_o\parallel\frac{r_\pi+(R_s\parallel R_B)}{\beta+1}=47.01\Oh.\]
The gain is positive and less than unity; the follower supplies impedance transformation. Because the collector is AC ground, the collector--emitter resistor $r_o$ is an emitter-to-ground resistor. Include it in the loaded emitter resistance for gain, and retain it after the external load is removed for output resistance.\par\vfill

\newpage
\head{Worked solutions: 23--24}
\sol{23}{Design for a specified loaded gain}$g_m=.04$ S, $r_\pi=2500\Oh$, and $R_C\parallel R_L=2000\Oh$. Use the exact finite-$\beta$ expression:
\[-10=-\frac{100(2000)}{2500+101R_E}\quad\Rightarrow\quad
R_E=\frac{20000-2500}{101}=173.3\Oh.\]
\[R_{\rm in}=50\K\parallel20\K=14290\Oh,\quad
A_s=-10\frac{14290}{1000+14290}=-9.346,\quad R_{\rm out}=4\K.\]
The source divider still reduces $v_i$ relative to $v_s$. The question fixes a loaded stage gain, not a source-to-output gain.\par\vfill
\sol{24}{Diagnose a complete wrong solution}First $g_m=40\mS$ and $r_\pi=2500\Oh$. The transistor input path is
\[R_{ib}=r_\pi+(\beta+1)R_E=12600\Oh,\quad R_{\rm in}=R_B\parallel R_{ib}=6848\Oh.\]
The effective collector load is $R_C\parallel R_L=2000\Oh$. Hence
\[A_v=-\frac{\beta(2000)}{2500+(\beta+1)(100)}=-15.87,\qquad
A_s=A_v\frac{6848}{1000+6848}=-13.85.\]
With the load removed and the source zeroed, $R_{\rm out}=3000\Oh$. Mistakes: (1) $R_s$ is outside the amplifier input port, $R_B$ is omitted, and emitter resistance is not reflected to the base. (2) Gain omits both the collector load and degeneration. (3) Source attenuation is omitted. (4) The external load must be removed for output resistance. The corrected output peak amplitude is $13.85$ mV for a 1 mV source.\par\vfill
\end{document}